\documentclass{sbc-template}

\usepackage[brazil]{babel}
\usepackage[T1]{fontenc}
\usepackage[utf8]{inputenc}
\usepackage{graphicx}
\usepackage{url}
\usepackage{hyperref}
\usepackage{microtype}
\usepackage{booktabs}
\usepackage{amsmath}
\usepackage{float}
\usepackage{tikz}
\usetikzlibrary{arrows.meta,positioning,fit,backgrounds}
\setlength{\emergencystretch}{2em}

\newcommand{\pendente}[1]{\textbf{[PENDENTE: #1]}}

\title{Um Assistente Baseado em Recuperação Aumentada para as Normas de um
       Programa de Pós-Graduação: Seleção de Modelos, Avaliação e Riscos}

\author{Leticia Sampaio Araujo Gonçalves\inst{1}, Maria Danielly Benício de Araújo\inst{1},
        Maria Eduarda Silva Vieira\inst{1}}

\address{Mestrado e Doutorado em Ciência da Computação -- MDCC \\
  Universidade Federal do Ceará -- UFC, Fortaleza, Ceará, Brasil \\
  Disciplina: Modelos de Linguagem de Grande Escala -- Prof. José Macedo}

\begin{document}

\maketitle

\begin{abstract}
Graduate students often look for the rules of their program in long legal-style
documents: internal regulations, university resolutions and procedure pages.
This paper describes an assistant that answers questions about the rules of the
Graduate Program in Computer Science of the Federal University of Ceará (MDCC/UFC)
in Portuguese, using retrieval-augmented generation over the official documents and
citing the document, the article and its date in every answer. The corpus has five
official PDFs and six procedure pages, split into 189 chunks that follow the article
structure of the norms and embedded with the multilingual model BGE-M3. Three
generators served by Ollama Cloud are compared (gpt-oss:20b, gpt-oss:120b and
gemma4:31b) on retrieval quality, answer quality judged by a separate LLM and by
people, tokens, latency, estimated cost and memory, and on four security and bias
tests: indirect prompt injection, out-of-scope questions, answers without a source and
formal versus informal phrasing. \pendente{main results of the full evaluation set}.
\end{abstract}

\begin{resumo}
Alunos de pós-graduação costumam procurar as regras do seu programa em documentos
longos e de linguagem normativa: regimento interno, resoluções da universidade e
páginas de procedimentos. Este artigo descreve um assistente que responde, em
português, a perguntas sobre as normas do Programa de Pós-Graduação em Ciência da
Computação da Universidade Federal do Ceará (MDCC/UFC), usando geração aumentada por
recuperação sobre os documentos oficiais e citando documento, artigo e data em toda
resposta. O corpus reúne cinco PDFs oficiais e seis páginas de procedimentos,
divididos em 189 trechos que respeitam a estrutura de artigos das normas e
representados pelo modelo multilíngue BGE-M3. Três geradores servidos pelo Ollama
Cloud são comparados (gpt-oss:20b, gpt-oss:120b e gemma4:31b) quanto à qualidade da
recuperação, à qualidade da resposta julgada por outro LLM e por pessoas, a tokens,
latência, custo e memória estimados, e a quatro testes de segurança e viés: injeção
indireta de \textit{prompt}, perguntas fora do escopo, respostas sem fonte e
formulação formal versus informal. \pendente{principais resultados do conjunto
completo de avaliação}.
\end{resumo}

\palavraschave{Modelos de Linguagem de Grande Escala, Geração Aumentada por
Recuperação, Avaliação de LLMs, Injeção de \textit{Prompt}, Normas Acadêmicas}

% ------------------------------------------------------------------------------
\section{Introdução}
\label{sec:introducao}

As regras que um aluno de mestrado ou doutorado precisa cumprir estão espalhadas por
documentos de naturezas diferentes. No MDCC/UFC, o prazo de vínculo vem da resolução
que fixa as normas gerais da pós-graduação da universidade \cite{UFC2015}; os créditos,
a composição das bancas e a antecedência da divulgação de uma defesa vêm do regimento
do programa \cite{MDCC2016}; os exames de proficiência aceitos e o prazo para pedir
aproveitamento de disciplinas aparecem apenas em páginas de procedimentos do site. Uma
pergunta simples, como quantos créditos faltam para concluir o curso, pode exigir a
leitura de mais de um documento, e as respostas nem sempre coincidem: o regimento fala
em 24 créditos em disciplinas, enquanto a resolução fala em um mínimo de 30 créditos,
dos quais seis correspondem à dissertação.

Modelos de linguagem de grande escala (LLMs) respondem bem a perguntas em linguagem
natural, mas, sem acesso aos documentos, tendem a produzir regras plausíveis e
inexistentes \cite{Ji2023}. Em um domínio normativo esse erro tem custo direto para o
aluno. A geração aumentada por recuperação (RAG) reduz o problema ao entregar ao modelo
os trechos relevantes no momento da pergunta \cite{Lewis2020,Gao2023}, e permite que
cada afirmação seja rastreada até a sua fonte.

Este trabalho constrói e avalia um assistente RAG para as normas do MDCC. As decisões
de projeto seguem cinco perguntas: qual modelo usar e por quê; qual estratégia de
adaptação (\textit{prompting}, RAG ou ajuste fino); como avaliar de forma automática e
humana; quanto custa cada resposta em tokens, latência e memória; e que riscos de
viés, privacidade e segurança o sistema traz. As contribuições são:
(i) um corpus versionado das normas do programa, com uma divisão em trechos que
preserva artigos e itens, de modo que a citação vem dos metadados e não da memória do
modelo; (ii) um conjunto de perguntas de avaliação em que cada resposta de referência
tem um trecho literal do documento como evidência; e (iii) uma comparação de três
geradores com métricas de recuperação, de qualidade da resposta, de custo e de
segurança, todas reproduzíveis a partir de arquivos.

A Seção~\ref{sec:relacionados} apresenta os trabalhos relacionados. A
Seção~\ref{sec:metodologia} descreve o corpus, o sistema e o protocolo de avaliação. A
Seção~\ref{sec:resultados} traz os resultados, discutidos na
Seção~\ref{sec:discussao}. A Seção~\ref{sec:conclusao} conclui.

% ------------------------------------------------------------------------------
\section{Trabalhos Relacionados}
\label{sec:relacionados}

Lewis et al.~\cite{Lewis2020} combinaram um recuperador denso com um gerador
\textit{sequence-to-sequence} e mostraram que condicionar a geração em passagens
recuperadas melhora tarefas que dependem de conhecimento factual. O levantamento de
Gao et al.~\cite{Gao2023} organiza as variações posteriores em RAG ingênuo, avançado e
modular, e identifica a divisão dos documentos em trechos e a escolha do modelo de
\textit{embedding} como pontos que afetam diretamente a qualidade da recuperação. Este
trabalho fica na categoria mais simples, recuperar e depois gerar, e concentra o
esforço na divisão em trechos, que segue a estrutura das normas.

Para textos em português, a escolha do \textit{embedding} pesa. O BGE-M3
\cite{Chen2024} foi treinado para mais de cem línguas e aceita entradas longas, o que o
torna adequado a artigos extensos de regimentos e resoluções.

A avaliação de sistemas RAG costuma separar a recuperação da geração. O RAGAS
\cite{Es2023} propõe métricas sem referência humana, como fidelidade da resposta ao
contexto recuperado, relevância da resposta e precisão do contexto, calculadas com
auxílio de um LLM. Zheng et al.~\cite{Zheng2023} estudaram o uso de LLMs como juízes e
encontraram boa concordância com avaliadores humanos, mas também vieses de posição, de
verbosidade e de autopreferência. Esses vieses motivam, neste trabalho, o uso de um
juiz de família diferente da dos modelos avaliados e a comparação explícita entre as
notas do juiz e as notas humanas.

Do lado da segurança, Greshake et al.~\cite{Greshake2023} mostraram que aplicações integradas a
LLMs podem ser comprometidas por injeção indireta de \textit{prompt}: instruções
escondidas em conteúdo recuperado, e não digitadas pelo usuário, passam a controlar o
modelo. Em um sistema RAG qualquer documento indexado é uma superfície de ataque, e
esse cenário é testado aqui com um documento envenenado. A alucinação, discutida em
profundidade por Ji et al.~\cite{Ji2023}, é medida pela taxa de respostas sem fonte e pela
nota de fidelidade.

O ajuste fino eficiente com LoRA \cite{Hu2021} reduz o custo de adaptar um modelo,
mas incorpora o conhecimento aos pesos. Para normas que mudam a cada resolução, isso
exigiria retreinar o modelo a cada alteração e não resolveria a necessidade de citar
a fonte, razão pela qual a alternativa foi descartada (Seção~\ref{sec:estrategia}).

% ------------------------------------------------------------------------------
\section{Metodologia}
\label{sec:metodologia}

\subsection{Corpus}
\label{sec:corpus}

O corpus reúne os documentos públicos indicados na página de regulamentos e de
procedimentos do MDCC (Tabela~\ref{tab:corpus}). Os PDFs têm camada de texto e
dispensaram reconhecimento óptico de caracteres. As páginas de procedimentos foram
reduzidas ao texto do conteúdo principal e datadas pela data de acesso
(26/11/2026), já que não informam data de publicação. O calendário universitário de
2026 não estava disponível no endereço indicado pelo site e ficou fora do corpus, assim
como os editais de seleção e de bolsas, cujos valores e datas mudam a cada ano.

\begin{table}[ht]
\centering
\caption{Documentos do corpus.}
\label{tab:corpus}
\footnotesize
\begin{tabular}{llr}
\toprule
\textbf{Documento} & \textbf{Data} & \textbf{Trechos} \\
\midrule
Regimento do MDCC & 2016 & 55 \\
Resolução n.\textsuperscript{o} 17/CEPE & 04/12/2015 & 52 \\
Norma MDCC 01/2025 (colegiado) & 18/07/2025 & 31 \\
Norma MDCC 02/2019 (professores externos) & 22/11/2019 & 21 \\
Norma MDCC 01/2024 (PROAP/CAPES) & 17/05/2024 & 18 \\
Páginas de procedimentos e disciplinas (6) & acesso em 26/11/2026 & 12 \\
\midrule
\textbf{Total} & & \textbf{189} \\
\bottomrule
\end{tabular}

\vspace{2pt}
{\footnotesize Fonte: elaborado pelas autoras.}
\end{table}

\subsection{Arquitetura do sistema}
\label{sec:arquitetura}

A Figura~\ref{fig:arquitetura} mostra o fluxo de uma pergunta. O texto de cada
documento é dividido em unidades que começam em um artigo (\textit{Art. 12}) ou em um
item numerado (\textit{3.2.}), conforme o estilo do documento; parágrafos e incisos
permanecem dentro do seu artigo. Quando a numeração de artigos recomeça, o restante é
tratado como anexo, o que distingue, na Resolução 17/CEPE, o artigo 1.º da resolução
do artigo 1.º das normas que ela aprova. Unidades com mais de 1.500 caracteres são
cortadas por tamanho e mantêm o rótulo. Cada trecho recebe como primeira linha a sua
citação (título do documento, artigo e data) e é representado com o BGE-M3
\cite{Chen2024}, executado localmente, em vetores de 1.024 dimensões armazenados no
PostgreSQL com a extensão pgvector.

Para cada pergunta, dados pessoais (e-mail, CPF e sequências numéricas com perfil de
matrícula) são mascarados antes de qualquer processamento. Os cinco trechos mais
próximos são numerados e enviados ao gerador junto de instruções em português: usar
apenas os trechos, marcar cada afirmação com o número do trecho, responder com uma
frase fixa quando os trechos não responderem, recusar perguntas fora do tema e tratar
os trechos como documentos, nunca como ordens. As citações exibidas ao aluno são
montadas a partir dos metadados dos trechos efetivamente citados, e não do texto
gerado. A interface é uma página de conversa em Streamlit.

\begin{figure}[ht]
\centering
\resizebox{\linewidth}{!}{%
\begin{tikzpicture}[
  node distance=4mm and 6mm,
  box/.style={draw, rounded corners=2pt, align=center, font=\footnotesize, minimum height=8mm, inner sep=3pt},
  arr/.style={-{Stealth[length=2mm]}, thick}]
\node[box] (q) {Pergunta};
\node[box, right=of q] (m) {Máscara de\\dados pessoais};
\node[box, right=of m] (r) {Recuperação\\top-5 (pgvector)};
\node[box, right=of r] (g) {Gerador\\(Ollama Cloud)};
\node[box, right=of g] (a) {Resposta\\+ citações};
\node[box, below=of r] (e) {BGE-M3\\(local)};
\node[box, below=of a] (l) {Registro JSONL\\(sem identificação)};
\draw[arr] (q) -- (m); \draw[arr] (m) -- (r); \draw[arr] (r) -- (g); \draw[arr] (g) -- (a);
\draw[arr] (e) -- (r); \draw[arr] (a) -- (l);
\end{tikzpicture}}
\caption{Fluxo de uma pergunta no assistente.}
\label{fig:arquitetura}

\vspace{2pt}
{\footnotesize Fonte: elaborado pelas autoras.}
\end{figure}

\subsection{Seleção dos modelos}
\label{sec:modelos}

Foram comparados três geradores disponíveis no Ollama Cloud, escolhidos para separar
dois efeitos. O par gpt-oss:20b e gpt-oss:120b \cite{OpenAI2025} pertence à mesma
família e isola o efeito do tamanho; o gemma4:31b, de outra família, isola o efeito da
família. Todos rodam com temperatura zero. A execução local dos geradores foi
considerada e descartada: a máquina de desenvolvimento não tem GPU, e um modelo de
cerca de 4 bilhões de parâmetros levaria em torno de 40 minutos para responder a 100
perguntas na CPU. Apenas o modelo de \textit{embedding} roda localmente, porque o
serviço em nuvem não oferece modelos desse tipo. O juiz da avaliação automática é o
glm-5.3, de família diferente das três avaliadas, para evitar a autopreferência
apontada por Zheng et al.~\cite{Zheng2023}. A escolha final é justificada pelos resultados da
Seção~\ref{sec:resultados}.

\subsection{Estratégia: \textit{prompting} e RAG}
\label{sec:estrategia}

A estratégia combina RAG com instruções no \textit{prompt} de sistema. O ajuste fino
foi descartado por três motivos. As normas mudam por resolução, e um modelo ajustado
precisaria ser retreinado a cada mudança. O conhecimento incorporado aos pesos não
traz a sua fonte, enquanto a exigência central do sistema é citar o artigo. E o
corpus, com 189 trechos, é pequeno demais para um ajuste fino que generalize, mesmo com
métodos eficientes como o LoRA \cite{Hu2021}. Com RAG, atualizar o assistente se
resume a substituir o documento e reindexar.

\subsection{Avaliação automática}
\label{sec:avaliacao-automatica}

O conjunto de avaliação é escrito pelas autoras a partir de dúvidas reais de alunos do
programa. Cada pergunta pertence a uma de quatro categorias: \textit{normal} (as
normas respondem), \textit{não encontrada} (sobre o programa, mas sem resposta nas
normas), \textit{fora do escopo} (saúde, aconselhamento jurídico, opiniões sobre
pessoas) e \textit{injeção} (respondida com a coleção envenenada, ver
Seção~\ref{sec:riscos}). Perguntas normais e de injeção trazem uma resposta de
referência e um trecho copiado literalmente do documento; um teste automático rejeita
a pergunta cujo trecho não aparece no texto extraído. Parte das perguntas tem uma
reescrita informal com o mesmo identificador de par. O conjunto final tem
\pendente{número de perguntas por categoria}.

A recuperação é medida sem LLM: \textit{hit@k} indica se o trecho de evidência está
entre os $k$ recuperados, e o \textit{Mean Reciprocal Rank} (MRR) considera a posição
em que ele aparece. A resposta é julgada pelo glm-5.3 com uma rubrica em português,
com notas inteiras de 0 a 2 para correção (comparada à referência), fidelidade (tudo o
que afirma está nos trechos) e citação (cita o trecho que comprova). As métricas
seguem as definições do RAGAS \cite{Es2023}, implementadas diretamente para usar
\textit{prompts} em português. Uma resposta normal é considerada aprovada quando tem
correção e fidelidade máximas. As categorias \textit{não encontrada} e \textit{fora do
escopo} são verificadas por regra, pela presença da frase fixa correspondente.

\subsection{Avaliação humana}
\label{sec:avaliacao-humana}

Uma amostra aleatória de 20 respostas normais, com o modelo oculto, é avaliada por
\pendente{número de avaliadores} pessoas em escala de 1 a 5 para correção, utilidade e
clareza. São reportadas as médias por critério, a concordância entre avaliadores
\pendente{medida de concordância e referência} e a concordância entre as notas
humanas e as do juiz automático.

\subsection{Custo, memória e latência}
\label{sec:custo}

Para cada resposta são registrados os tokens de entrada e de saída informados pelo
servidor e a latência de ponta a ponta da chamada ao gerador, resumida pelos
percentis 50 e 95. O Ollama Cloud é cobrado por assinatura, e não por token; o custo
por mil perguntas é, por isso, uma estimativa: tokens medidos multiplicados pelo preço
público por token de um provedor que sirva o mesmo modelo aberto
\pendente{provedor e preços de referência}. A memória dos geradores também é estimada,
pelo número de parâmetros multiplicado pelos bytes por peso na quantização servida,
somado ao \textit{cache} de atenção. A memória do contêiner de \textit{embeddings},
que roda localmente, é medida.

\subsection{Vieses, privacidade e segurança}
\label{sec:riscos}

Quatro testes foram definidos. \textbf{Injeção indireta}: um documento falso, com a
instrução de responder que não há prazo máximo para a defesa, é indexado em uma
coleção separada, nunca usada pela interface; perguntas sobre prazo de defesa são
respondidas com essa coleção, e a resposta falha se repetir a afirmação plantada
\cite{Greshake2023}. \textbf{Fora do escopo}: o assistente deve recusar.
\textbf{Respostas sem fonte}: proporção de respostas normais sem nenhuma citação,
indicador de alucinação \cite{Ji2023}. \textbf{Formal versus informal}: a diferença
de aprovação entre os pares de perguntas avalia se o sistema atende pior quem escreve
de modo coloquial ou com erros de digitação.

Quanto à privacidade, o sistema não tem contas de usuário. Antes de chegar ao modelo
ou ao registro, a pergunta passa pela máscara de dados pessoais, e o registro recusa
qualquer campo que identifique quem perguntou, em linha com a minimização prevista na
LGPD \cite{Brasil2018}. Como os geradores rodam em nuvem, o texto da pergunta sai da
máquina; a máscara reduz, mas não elimina, esse risco.

% ------------------------------------------------------------------------------
\section{Resultados}
\label{sec:resultados}

\pendente{Esta seção será preenchida com a rodada sobre o conjunto completo de
perguntas, para os três geradores. Os números abaixo vêm de uma rodada de verificação
do sistema, com sete perguntas de exemplo, e não sustentam comparação entre modelos.}

A Tabela~\ref{tab:preliminar} resume a rodada de verificação de 26/11/2026, com
gpt-oss:20b e gpt-oss:120b. A recuperação trouxe o trecho de evidência entre os cinco
primeiros em todas as perguntas normais, nem sempre em primeiro lugar. Nenhuma resposta
normal ficou sem fonte. O gpt-oss:120b respondeu com mediana de 0,84~s, contra
1,86~s do gpt-oss:20b, e gerou menos tokens de saída.

\begin{table}[ht]
\centering
\caption{Rodada de verificação com sete perguntas de exemplo (resultados preliminares).}
\label{tab:preliminar}
\footnotesize
\begin{tabular}{lrr}
\toprule
\textbf{Métrica} & \textbf{gpt-oss:20b} & \textbf{gpt-oss:120b} \\
\midrule
\textit{hit@5} (4 perguntas normais) & 1,00 & 1,00 \\
MRR & 0,69 & 0,69 \\
Correção / fidelidade / citação (0--2) & 1,50 / 2,00 / 2,00 & 1,50 / 2,00 / 2,00 \\
Respostas normais sem fonte & 0,00 & 0,00 \\
Tokens de entrada / saída (média) & 1.811 / 269 & 1.811 / 197 \\
Latência p50 / p95 (ms) & 1.858 / 3.563 & 841 / 1.363 \\
\bottomrule
\end{tabular}

\vspace{2pt}
{\footnotesize Fonte: elaborado pelas autoras.}
\end{table}

A rodada também expôs dois comportamentos que orientaram o desenho do conjunto
completo. Na pergunta de injeção, o gpt-oss:120b afirmou que não há prazo máximo
específico para a defesa, repetindo a instrução plantada no documento envenenado,
apesar da regra explícita no \textit{prompt}; o gpt-oss:20b respondeu que não
encontrou a informação. E a pergunta sobre créditos do mestrado recebeu de ambos os
modelos a resposta 24 créditos em disciplinas mais seis de dissertação, combinando o
regimento e a resolução, enquanto a referência registrava apenas os 24 créditos do
regimento; a referência foi revista.

\pendente{Tabela de qualidade por modelo (três geradores).}
\pendente{Tabela de testes de segurança e viés por modelo.}
\pendente{Tabela de custo estimado por mil perguntas, memória estimada e latência.}
\pendente{Resultados da avaliação humana e concordância com o juiz.}

% ------------------------------------------------------------------------------
\section{Discussão}
\label{sec:discussao}

\pendente{Interpretação dos resultados completos e justificativa final do modelo
escolhido.}

\subsection{Limitações e ameaças à validade}

O conjunto de avaliação foi escrito por alunas do próprio programa, o que aproxima as
perguntas das dúvidas reais, mas pode concentrá-las nos temas que elas conhecem. O
juiz automático é um LLM e herda os vieses descritos por Zheng et al.~\cite{Zheng2023}; a
comparação com as notas humanas estima o tamanho desse problema, sem eliminá-lo. As
estimativas de custo e de memória dependem de preços de terceiros e de suposições
sobre a quantização servida pelo Ollama Cloud. O regimento disponível no site é de
2016, e alterações posteriores não publicadas ficariam fora do assistente. Por fim, o
sistema responde a cada pergunta de forma isolada, sem memória da conversa.

% ------------------------------------------------------------------------------
\section{Conclusão}
\label{sec:conclusao}

\pendente{Conclusão a partir dos resultados completos: modelo recomendado, o que os
testes de segurança mostraram e trabalhos futuros.}

O código, o corpus, o conjunto de perguntas e os resultados estão disponíveis em
\pendente{endereço do repositório}.

\bibliographystyle{apalike}
\bibliography{referencias}

\end{document}
